\documentclass[aspectratio=169,10pt]{beamer}
\usetheme{metropolis}
\metroset{block=fill, sectionpage=progressbar}
\setbeamersize{text margin left=0.7cm, text margin right=0.7cm}
\usepackage{fontspec}
\setmainfont{Noto Serif CJK KR}
\setsansfont{Noto Sans CJK KR}
\setmonofont{DejaVu Sans Mono}
\usepackage{amsmath, amssymb}
\usepackage{xcolor}
\usepackage{hyperref}

\definecolor{aprlblue}{HTML}{1F4E79}
\definecolor{aprllight}{HTML}{4A7FB5}
\definecolor{aprlaccent}{HTML}{D9531E}
\definecolor{aprlgray}{HTML}{4A4A4A}
\definecolor{aprlbg}{HTML}{F2F4F7}

\setbeamercolor{frametitle}{bg=aprlblue, fg=white}
\setbeamercolor{title separator}{fg=aprlaccent}
\setbeamercolor{progress bar}{fg=aprlaccent, bg=aprlbg}
\setbeamercolor{alerted text}{fg=aprlaccent}

% 대본 본문용 환경 - 읽기 좋게
\newenvironment{script}{%
  \footnotesize
  \setlength{\parskip}{0.5em}
  \setlength{\parindent}{0pt}
}{}

\hypersetup{colorlinks=true, urlcolor=aprllight}

\title{VLA 강의 --- 발표 대본}
\subtitle{RT2에서 $\pi_0$까지, 그리고 그 너머}
\author{Giseop Kim}
\institute{APRL, DGIST \quad $\cdot$ \quad RT604 Spring 2026}
\date{}

\begin{document}

% ============================================================
{
\setbeamertemplate{footline}{}
\begin{frame}[plain]
\vspace*{0.8cm}
\begin{center}
{\LARGE\textbf{VLA 강의 --- 발표 대본}}\\[0.3cm]
{\large RT2에서 $\pi_0$까지, 그리고 그 너머}\\[0.8cm]
\textbf{Giseop Kim}\\[0.15cm]
{\small APRL, DGIST \textperiodcentered{} RT604 Spring 2026}
\end{center}

\vspace{1.0cm}
\begin{block}{출처 (Source)}
본 강의 자료는 다음 영상을 기반으로 재구성되었습니다.
\begin{itemize}
  \footnotesize
  \item \textbf{제목}: \textit{Inside the World's Smartest Robot Brain [VLA]}
  \item \textbf{저자}: \href{https://www.youtube.com/@WelchLabs}{Welch Labs}
  \item \textbf{영상 URL}: \href{https://www.youtube.com/watch?v=2mrGMMmrVNE}{\texttt{youtube.com/watch?v=2mrGMMmrVNE}}
\end{itemize}
\end{block}
\end{frame}
}

% ============================================================
\begin{frame}{이 문서에 대하여}
\begin{itemize}
  \item 이 문서는 \texttt{vla\_lecture.tex}(슬라이드 본편)와 \textbf{1:1로 대응}되는 발표 대본입니다.
  \item 한 슬라이드당 \textbf{2--4분 분량}의 구어체 설명을 담았습니다 (총 약 60--90분 강의 분량).
  \item 강의 시 \texttt{vla\_lecture.pdf}를 띄워놓고 이 문서를 \textbf{진행 노트}로 활용하시면 됩니다.
  \item 대본은 그대로 읽기보다는 \alert{흐름과 강조점을 잡는 가이드}로 사용하기를 권장합니다.
\end{itemize}

\vspace{1em}
\textbf{표기}
\begin{itemize}
  \item 각 슬라이드 대본 시작 부분에 \textbf{[\,슬라이드 번호 / 제목\,]} 표시
  \item \textit{기울임체}: 청중의 반응을 유도하는 질문 또는 강조 톤
  \item [말줄임]: 호흡을 두는 자리
\end{itemize}
\end{frame}

% ============================================================
\section{Part 1. Coke Can과 Taylor Swift}

\begin{frame}{[3] 2023년 7월: 어색했지만 결정적인 데모}
\begin{script}
자, 본격적으로 들어가 봅시다. 시계를 2023년 7월로 돌려볼게요.

Google의 한 연구실, 평범한 책상 위에 코크 캔 하나, 그리고 Tom Cruise, Snoop Dogg, Taylor Swift의 사진 세 장이 놓여 있습니다. 연구자가 로봇에게 명령을 내립니다. \textit{``Move the Coke can to Taylor Swift.''}

이 명령을 받은 로봇은 RT2라는 새로운 두뇌로 움직였는데요, 흥미로운 점은 RT2가 너무 커서 \alert{로봇 안에 올릴 수가 없었다}는 겁니다. 그래서 카메라 영상을 한 장씩 클라우드의 TPU 클러스터로 보내고, 거기서 제어 신호를 받아오는 방식이었어요.

그렇게 천천히, 정말 어색하게, 로봇이 코크 캔을 들어서 Taylor Swift 사진의 가장자리에 슬쩍 내려놓았습니다.

[말줄임] 솔직히 말하면 영상을 보면 시시해 보입니다. 그런데 2년 후, 당시 팀 멤버였던 Carol Houseman이 인터뷰에서 이 장면을 이렇게 회고합니다. \textit{``이게 정말 될 거라는 확신이 든 순간이었어요.''}

그 후 1년 안에 RT2 핵심 멤버 대부분이 Google을 떠나서 \textbf{Physical Intelligence}라는 스타트업을 차리게 됩니다. 이 어색한 데모가 어떻게 한 회사를 탄생시킬 만한 사건이었을까요? 다음 슬라이드에서 그 답을 봅시다.
\end{script}
\end{frame}

% ============================================================
\begin{frame}{[4] 왜 이 어색한 데모가 돌파구였나}
\begin{script}
핵심은 한 가지 사실에 있습니다. \alert{로봇 학습 데이터에 Taylor Swift는 단 한 번도 등장하지 않았다}는 것.

생각해 보세요. RT2가 이 명령을 수행하려면 두 종류의 지식을 동시에 갖고 있어야 합니다.

\textbf{첫째}, ``Taylor Swift는 사람이고, 이 얼굴이 Taylor Swift다''라는 의미적 지식. 이건 인터넷 사전학습에서 옵니다. 위키피디아, 뉴스 기사, 이미지에 달린 캡션, 이런 데이터요.

\textbf{둘째}, ``코크 캔을 집어서 어떤 위치로 옮긴다''는 물리적 행동. 이건 인간이 로봇을 조작한 시연 데이터에서 옵니다.

이 둘은 \textbf{전혀 다른 종류의 데이터}예요. 한쪽은 텍스트와 이미지, 다른 쪽은 관절 각도와 그리퍼 명령. RT2가 이 작업을 해냈다는 건, 모델 안에서 이 두 세계가 \alert{연결}되었다는 뜻입니다.

[말줄임] 이게 왜 중요한가? 만약 정말로 이게 가능하다면, 우리가 가진 \textbf{인터넷 전체의 지식}을 로봇 두뇌에 부어 넣을 수 있다는 얘기거든요. 학습 데이터를 일일이 수집할 필요가 없어진다는 의미죠. 위키피디아에 있는 모든 개념이 로봇이 활용할 수 있는 지식이 되는 겁니다.

대학원생 입장에서 생각해 봅시다. 여러분이 로봇 학습 데이터셋을 만든다고 하면 가장 큰 병목이 뭐죠? \textit{데이터 수집입니다.} VLA는 그 병목을 우회할 가능성을 보여준 첫 번째 정량적 증거였어요.
\end{script}
\end{frame}

% ============================================================
\section{Part 2. Google의 진화 경로}

\begin{frame}{[5] 2022: SayCan --- LLM을 플래너로}
\begin{script}
이제 RT2가 어떤 흐름에서 나왔는지를 봅시다. 출발점은 2022년의 SayCan이에요.

ChatGPT가 막 나오던 시기죠. Google 연구자들이 자연스럽게 떠올린 질문이 있었어요. \textit{``LLM이 이렇게 똑똑한데, 로봇 플래닝에 쓸 수 없을까?''}

SayCan의 아이디어는 간단합니다. 사람이 ``엎질러진 것을 닦아라''라고 말하면, LLM이 이걸 단계별로 쪼개주는 거예요. ``스펀지를 찾아라, 스펀지를 집어라, 엎질러진 곳으로 가라, 닦아라.'' 이렇게요.

근데 여기 결정적인 한계가 있었습니다. \alert{실제 제어는 별도의 imitation learning 네트워크가 했다}는 점이에요. LLM은 ``스펀지를 집어라''라고 말만 할 수 있고, 실제로 그게 가능하려면 별도의 네트워크가 ``스펀지 집기''를 미리 학습해 둬야 했습니다.

즉, LLM은 \textbf{미리 정해진 메뉴판}에서 항목을 고르는 역할일 뿐이었어요. 메뉴에 없는 행동은 못 합니다. 코크 캔을 Taylor Swift 사진에 올리려면? 그 행동을 미리 학습시켜야 했어요.

[말줄임] 또 하나 중요한 한계는, SayCan에서 \textbf{LLM이 눈을 감고 있었다}는 거예요. 텍스트 전용 모델이라 카메라 이미지를 볼 수 없었거든요. 이게 나중에 PaLM-E로 풀리는 문제입니다.
\end{script}
\end{frame}

% ============================================================
\begin{frame}{[6] 2022 말: RT1 --- 더 큰 행동 메뉴}
\begin{script}
SayCan의 메뉴 문제를 해결하려면 어떻게 해야 할까요? 답은 단순합니다. \textbf{메뉴를 키운다.}

2022년 말 Google이 발표한 RT1, Robot Transformer 1이 정확히 그 일을 했어요. 같은 imitation learning 패러다임이지만, 데이터를 \alert{13만 건 이상}으로 대규모로 모았습니다. 그리고 transformer 기반 아키텍처를 도입했죠.

이게 SayCan의 플래너 LLM에게 훨씬 풍부한 ``할 줄 아는 것들의 목록''을 제공해 주었습니다. 처음 보는 주방에서 특정 물건을 찾는 작업처럼, 여러 단계가 필요한 long-horizon task에서 성능이 눈에 띄게 좋아졌어요.

[말줄임] 그런데도 한 가지 문제가 여전히 남아 있었어요. 플래너 LLM은 \textbf{여전히 텍스트만 봅니다}. 시각 정보가 없어요. ``서랍을 열어라''라고 계획을 세웠는데 누가 이미 서랍을 열어놨다? LLM은 모릅니다. 계속 ``서랍을 열어라''라고 지시할 거예요.

이게 다음 단계인 PaLM-E의 동기가 됩니다.
\end{script}
\end{frame}

% ============================================================
\begin{frame}{[7] 2023년 3월: PaLM-E --- 시각을 가진 플래너}
\begin{script}
2023년 3월 6일. GPT-4 출시 약 일주일 전이에요. Google이 PaLM-E를 발표합니다.

PaLM-E는 한 줄로 요약하면, \textbf{PaLM에 이미지 입력을 직접 통합한 멀티모달 LLM}이에요. 이제 플래너가 카메라를 통해 세상을 볼 수 있게 된 겁니다.

이게 얼마나 큰 차이를 만들었는지, 데모 하나를 소개할게요. 로봇에게 ``칩 봉지를 가져와라''라고 지시합니다. 로봇이 서랍을 열고 칩 봉지를 꺼내려는 순간, 연구자가 장난스럽게 봉지를 다시 서랍에 집어넣어요.

기존 시스템이라면 로봇은 자기가 이미 ``꺼냈다''고 생각하고 다음 단계로 넘어가겠죠. 그런데 PaLM-E는 \alert{상황이 바뀌었음을 인식하고 계획을 다시 세웁니다}. 다시 서랍을 열고, 다시 봉지를 잡고. 그것도 자율적으로요.

[말줄임] 이게 \textbf{visual feedback이 closed loop으로 들어온} 첫 번째 사례에 가깝습니다. 플래너가 매 순간 ``지금 세상이 어떤가?''를 보면서 계획을 적응시키는 거죠.

스택은 여전히 두 단계예요. PaLM-E가 플래닝하고, RT1이 제어합니다. 하지만 두 단계 모두 비전이 들어왔다는 게 큰 변화입니다.
\end{script}
\end{frame}

% ============================================================
\begin{frame}{[8] 관찰: 두 모델이 너무 닮았다}
\begin{script}
여기서 흥미로운 관찰이 시작됩니다. \textit{잠깐, PaLM-E랑 RT1이 너무 비슷하지 않나?}

두 모델 다 비전 인코더가 있고, transformer가 있고, 입력으로 이미지와 언어를 받아요. 차이는 단 하나, 출력입니다. PaLM-E는 텍스트를 출력하고, RT1은 제어 신호를 출력합니다.

자, 그럼 자연스러운 질문이 나오죠. \alert{``하나의 모델이 둘 다 출력하면 안 될까?''}

생각해 보세요. transformer의 출력은 결국 토큰 시퀀스예요. 텍스트 토큰을 출력하는 모델에, 행동을 표현하는 \textbf{action token}을 vocabulary에 추가하면 어떨까요? ``책상으로 가라''고 텍스트로 출력할 수도 있고, 관절 각도를 정수로 인코딩한 action token을 출력할 수도 있는, 그런 통합 모델을 만든다면?

[말줄임] 이게 단순히 엔지니어링 트릭이 아닙니다. 더 근본적인 질문이에요. \textit{지금까지 우리가 만든 가장 강력한 AI인 LLM이, 그 자체로 로봇 두뇌가 될 수 있는가?}

이 질문에 답을 시도한 게 다음 단계, RT2입니다.
\end{script}
\end{frame}

% ============================================================
\begin{frame}{[9] 2023년 7월: RT2 --- 통합의 탄생}
\begin{script}
2023년 7월. Google 로보틱스 팀이 RT2를 공개합니다. 레시피는 의외로 직설적이에요.

\textbf{1단계}: 멀티모달 LLM에서 출발한다. 구체적으로는 PaLM-E와 PaLI-X 두 가지를 베이스로 썼어요.

\textbf{2단계}: 6개월 전 RT1 학습에 사용했던 \emph{똑같은 인간 시연 데이터}를 가져온다.

\textbf{3단계}: 단, 출력을 토큰화된 제어 신호로 변환해서 fine-tune한다. 즉, ``오른팔을 5도 회전하라''를 하나의 토큰으로 표현하고, LLM이 이걸 텍스트 토큰처럼 출력하도록 학습시킨 거죠.

이게 됩니다. 그것도 \alert{학습 데이터에 없던 객체, 환경, 작업까지 일반화}하면서요. 우리가 처음 본 Taylor Swift 데모가 바로 그 결과입니다.

[말줄임] 팀은 이 새로운 종류의 모델에 이름을 붙입니다. \textbf{Vision-Language-Action}, 줄여서 VLA. 비전, 언어, 행동이 단일 모델로 통합됐다는 뜻이에요.

한 가지 강조할 게 있습니다. RT2 모델군은 50억에서 550억 파라미터 규모였어요. \textbf{로봇에 직접 올릴 수 없는 크기}입니다. 그래서 영상에서도 클라우드 추론을 했죠. 이게 다음 단계인 $\pi_0$의 또 다른 출발점이 됩니다.
\end{script}
\end{frame}

% ============================================================
\section{Part 3. $\pi_0$ 아키텍처 해부}

\begin{frame}{[10] 2024년 10월: $\pi_0$ --- 더 작고 더 빠르고 더 잘하는}
\begin{script}
RT2 데모로부터 15개월이 지난 2024년 10월. Physical Intelligence 팀이 $\pi_0$, ``Pi Zero''를 공개합니다.

차이가 인상적이에요. RT2는 코크 캔을 사진 위에 올리는 정도였다면, $\pi_0$는 \textbf{빨래를 건조기에서 꺼내 개고, 식탁을 정리하고, 그릴드 치즈 샌드위치를 만들고, 처음 보는 침실까지 정리}합니다. 정말로 ``집안일을 도와주는 로봇''에 가까워졌어요.

자, 그럼 $\pi_0$가 RT2보다 \textbf{더 큰 모델}일까요? 직관적으로는 그렇게 생각하기 쉽죠.

[말줄임] 정답은 반대입니다. RT2가 50억--550억 파라미터였는데, $\pi_0$는 \alert{33억 파라미터}예요. 더 작아졌습니다. 그리고 클라우드 TPU가 아니라 \alert{컨슈머급 RTX 4090} 한 장으로 \alert{73 ms}만에 추론합니다. 로봇 위에 직접 올라간다는 거죠.

\textit{어떻게 더 작은 모델이 더 잘할 수 있을까요?} 답은 ``아키텍처가 더 영리해졌다''에 있습니다. 다음 슬라이드에서 그 영리함을 봅시다.
\end{script}
\end{frame}

% ============================================================
\begin{frame}{[11] $\pi_0$ 전체 구조}
\begin{script}
$\pi_0$ 아키텍처를 한눈에 봅시다. 입력은 네 종류예요.

\textbf{왼쪽 상단}: 카메라 세 대(머리 위 + 양 손목)와 텍스트 프롬프트. 이건 \textbf{PaliGemma}로 들어갑니다. PaliGemma는 Google이 공개한 오픈 가중치 멀티모달 LLM이에요. SigLIP 이미지 인코더와 Gemma LLM이 결합된 형태죠.

\textbf{왼쪽 하단}: 로봇의 현재 상태(관절 각도 14개)와 노이즈 섞인 미래 행동들. 이건 \textbf{Action Expert}로 들어갑니다.

\textbf{오른쪽}: Action Expert가 14개 관절 × 50 시간 단계의 궤적을 출력합니다.

여기서 \alert{핵심은 두 모듈을 잇는 점선 화살표}입니다. ``KV cache''라고 적혀 있는 부분이요. PaliGemma가 계산한 정보가 Action Expert로 직접 전달되는 통로입니다. 이걸 어떻게 하는지가 $\pi_0$의 영리함의 중심이에요.

[말줄임] 슬라이드 아래의 세 가지 설계 결정을 강조하고 싶어요. 첫째, LLM 백본과 Action Expert를 분리. 둘째, flow matching으로 행동 생성. 셋째, KV cache로 두 모듈 결합. 다음 슬라이드들에서 하나씩 풀어봅시다.
\end{script}
\end{frame}

% ============================================================
\begin{frame}{[12] 설계 결정 1: 왜 Action Expert를 따로 두는가}
\begin{script}
첫 번째 설계 결정. \textit{왜 LLM이 직접 action token을 출력하지 않고, 별도의 Action Expert를 두는가?}

비유로 시작해 봅시다. 인간 뇌에서 ``전두엽의 인지''와 ``운동 피질의 운동 제어''는 다른 영역이 담당하잖아요. 분리되어 있지만 긴밀하게 협업합니다. $\pi_0$도 비슷한 철학이에요. \textbf{Gemma는 ``인지''에, Action Expert는 ``운동''에 특화}시킨 거죠.

흥미로운 디테일 하나 말씀드릴게요. Action Expert의 정체가 뭐냐 하면---\alert{Gemma와 똑같은 아키텍처}예요. 실제 코드에서도 Gemma 클래스로 인스턴스화됩니다. 차이는 두 가지뿐이에요. \textbf{무작위로 초기화된다}는 점, 그리고 \textbf{각 layer의 width가 더 작다}는 점입니다. 사전학습이 없으니 처음부터 행동에 맞게 학습되고, 작아서 가벼워서 실시간에 유리하죠.

[말줄임] 여기서 ``어, 이거 SayCan으로 돌아간 거 아니야?'' 하실 수 있어요. 둘 다 상위 모듈이 ``계획''하고 하위 모듈이 ``제어''하니까요. \textbf{핵심적인 차이는 인터페이스에 있습니다.}

SayCan에서는 두 모듈이 \alert{자연어로 통신}했어요. ``스펀지를 집어라''라는 텍스트로요. 이건 \textbf{매우 좁은 통로}입니다. 정보 손실이 커요.

$\pi_0$에서는 두 모듈이 \alert{attention 내부 표현으로 통신}합니다. 고차원 벡터 전체가 그대로 흐르는 거예요. 이게 ``하나의 두뇌처럼 사고하면서도 모듈성을 유지한다''는 말의 의미입니다.
\end{script}
\end{frame}

% ============================================================
\section{Part 4. 내부 동작: Attention head}

\begin{frame}{[13] Gemma LLM이 받는 입력}
\begin{script}
이제 모델 내부로 들어갑시다. Gemma LLM이 무엇을 입력으로 받는지부터요.

\textbf{이미지 처리}부터 봅시다. 카메라가 세 대 있어요---머리 위 1대, 양쪽 손목 1대씩. 각 이미지는 패치 그리드로 쪼개져서 이미지당 256개 패치가 됩니다. 세 장이니까 총 \textbf{768 패치}예요.

각 패치는 이미지 인코더(SigLIP)를 거쳐서 길이 2048짜리 임베딩 벡터로 변환됩니다. 이걸 \textbf{soft token}이라고도 불러요. ``부드럽게 인코딩된 토큰'' 정도의 의미죠.

\textbf{텍스트 처리}는 익숙하실 거예요. ``uncap the pen'' 같은 프롬프트가 4개 토큰으로 쪼개지고, 각 토큰이 임베딩 벡터로 매핑됩니다.

다 합치면 768 + 4 = \alert{772개의 soft token}이 Gemma로 들어가는 거예요.

[말줄임] Gemma 자체는 표준 transformer입니다. 18개 블록이 쌓여 있고, 각 블록은 attention과 MLP로 구성돼요. 각 attention 블록 안에는 8개의 head가 있습니다. 다음 두 슬라이드에서 이 head 하나가 무엇을 하는지 자세히 들여다볼 거예요.
\end{script}
\end{frame}

% ============================================================
\begin{frame}{[14] Attention head 한 개의 연산}
\begin{script}
Attention 한 head가 하는 연산을 수식으로 정리해 봅시다. 익숙한 분도 많겠지만, 다음 슬라이드의 ``pen'' 사례를 이해하기 위해 짚고 갑니다.

입력: 772개의 임베딩 벡터를 행으로 쌓은 행렬 $X$. 이게 출발점이에요.

\textbf{1단계 --- 세 행렬 생성}. $X$에 세 개의 학습 가능한 가중치 행렬을 곱해서 query $Q$, key $K$, value $V$ 세 행렬을 만듭니다. 각각 772행이에요. 입력 토큰 하나당 query 한 개, key 한 개, value 한 개씩이죠.

\textbf{2단계 --- Attention pattern}. $Q$와 $K^\top$를 곱하고 $\sqrt{d_k}$로 나눈 뒤 softmax를 씌웁니다. 결과는 772 × 772 행렬이에요. $i$번째 행, $j$번째 열의 값은 ``$i$번째 토큰이 $j$번째 토큰에 얼마나 주목하는가''를 나타냅니다.

\textbf{3단계 --- 출력}. attention pattern과 $V$를 곱합니다. 직관적으로는 \textit{각 query가 ``나에게 관련 있는 정보가 어디 있나?''를 묻고, 그 위치의 value를 자기 자리에 가져오는} 연산이에요.

[말줄임] 한 가지 더 강조할게요. $Q$와 $K$의 dot product가 큰 곳일수록 정보가 많이 ``복사''되는 겁니다. 다음 슬라이드에서 이 직관이 어떻게 ``펜 찾기''로 구현되는지 봅시다.
\end{script}
\end{frame}

% ============================================================
\begin{frame}{[15] 사례 연구: ``pen'' 토큰이 펜을 찾는다}
\begin{script}
이제 진짜 흥미로운 부분이에요. 학습된 attention head 하나를 \textbf{현미경으로 들여다본 결과}입니다.

시나리오: 프롬프트가 ``uncap the pen''이고, 마지막 토큰 ``pen''에 대응하는 query 벡터를 $q_\text{pen}$이라고 합시다.

이 head에서 무슨 일이 일어나느냐. $q_\text{pen}$과 모든 768개 image patch의 key 벡터 사이 dot product를 계산합니다. 결과를 보면 \alert{거의 모든 패치는 작은 값}이고, \alert{딱 두 패치에서만 값이 압도적으로 큽니다}. 그 두 패치가 어디냐? \textbf{실제로 펜이 보이는 그 두 패치}예요.

Softmax를 거쳐서 attention 값을 이미지 위에 히트맵으로 입혀 보면---펜이 있는 자리만 빨갛게 빛납니다. 다른 곳은 거의 검은색이고요.

여기서 \alert{진짜 놀라운 점}이 나와요. 이 분석을 비디오의 \emph{모든 프레임}에 대해 실행합니다. 그러면 어떻게 되겠어요? \textbf{펜이 자동으로 추적됩니다}. 별도의 추적 알고리즘 없이요. 그냥 attention head 하나가 그 일을 하는 거예요.

[말줄임] 이게 의미하는 바가 뭐냐. 이 모델은 \textbf{명시적으로 ``펜을 찾아라''라고 학습된 적이 없어요.} 그냥 인간 시연 데이터로 행동을 학습했을 뿐인데, 내부적으로 ``프롬프트의 명사를 이미지에서 grounding하는 기능''이 \alert{창발}한 겁니다.

마지막으로, 이건 단 하나의 head 이야기예요. 18 layer × 8 head = 144개의 head가 각자 다른 패턴을 학습하고 있습니다.
\end{script}
\end{frame}

% ============================================================
\begin{frame}{[16] Attention head의 의미론적 해석}
\begin{script}
앞 슬라이드에서 ``attention 값이 크면 정보가 복사된다''고 말씀드렸어요. 이걸 정확히 짚고 넘어갑시다.

수식으로 보면 출력 위치 $i$의 값에 $a_{ij}$ 곱하기 $v_j$가 더해집니다. 즉, ``pen'' 토큰 위치 $i$에 \textbf{펜 patch}의 value 벡터 $v_j$가 \alert{복사되어 합쳐지는} 거예요.

이걸 의미적으로 해석하면---\textbf{이 attention head는 텍스트의 ``pen''이라는 단어와, 이미지 속 펜이라는 시각적 객체를 하나의 통합 표현으로 묶어내고 있는} 겁니다. 언어와 시각의 grounding이 모델 내부 어딘가에서 이런 식으로 일어나고 있다는 거예요.

[말줄임] 그리고 다시 강조하지만, 우리가 본 건 \textbf{144개 중 단 하나의 head}예요. 다른 head들은요? 어떤 head는 ``손과 물체의 거리''에 주목할 수 있고, 어떤 head는 ``현재 작업의 진행 단계''를 추적할 수 있고, 어떤 head는 ``시간적 일관성''을 담당할 수 있어요.

연구자들이 이런 ``mechanistic interpretability'' 작업을 통해 LLM 내부에서 무슨 일이 벌어지는지를 조금씩 이해해 가고 있는데, VLA 모델에서도 이런 분석이 점점 더 중요해질 겁니다. 우리 lab에서 이런 방향으로 연구해 보고 싶은 분 있으면 따로 얘기해 봐요.
\end{script}
\end{frame}

% ============================================================
\section{Part 5. Action Expert와 Flow Matching}

\begin{frame}{[17] Action Expert의 입력}
\begin{script}
이제 Action Expert로 넘어갑니다. 입력이 두 종류예요.

\textbf{입력 1: 로봇 상태 (1 token)}. Aloha 플랫폼에서 각 팔은 7개 관절(허리, 어깨, 팔꿈치, 전완 회전, 손목, 손목 회전, 그리퍼)을 가져요. 양 팔 합치면 \textbf{14차원 벡터}입니다. 이 14차원 벡터를 14 × 1024 가중치 행렬로 곱해서, 길이 \alert{1024}짜리 임베딩 벡터 하나로 변환합니다.

여기서 짚고 넘어갈 점. Gemma의 임베딩은 길이 2048인데, Action Expert는 1024예요. \textbf{차원이 절반}이죠. 왜? \textbf{연산량과 추론 시간을 줄이기 위해서}입니다. 작은 모델로 충분하다는 판단인 거예요.

\textbf{입력 2: 노이즈가 섞인 미래 행동 (50 tokens)}. 다음 50 시간 단계의 14차원 관절 위치 예측이 50개 토큰으로 들어갑니다. 처음에는 \alert{순수한 노이즈}예요.

[말줄임] 잠깐, ``예측할 행동을 입력으로 받는다''는 게 직관적으로 이상하시죠? 이게 다음 슬라이드의 핵심 주제, flow matching입니다. 미리 예고하면, 이건 \textbf{이미지 생성 모델}에서 가져온 아이디어예요.
\end{script}
\end{frame}

% ============================================================
\begin{frame}{[18] Flow Matching: 노이즈에서 궤적으로}
\begin{script}
Flow matching의 직관은 단순합니다. 이미지 생성에서 \textit{``노이즈에서 시작해 점점 다듬으면 고양이 사진이 나온다''}는 그 메커니즘 있잖아요. \textbf{$\pi_0$는 같은 메커니즘으로 ``노이즈에서 시작해 점점 다듬으면 14×50 관절 궤적이 나온다''를 합니다.}

왜 이 방법이 잘 작동하느냐? 핵심은 \alert{분포가 multimodal}이라는 사실에 있어요.

이미지 생성에서 ``고양이 사진''이라는 분포를 생각해 보세요. 한 가지 ``올바른 고양이 이미지''가 아니라, 무수히 많은 가능한 고양이 이미지가 존재합니다. 정면 고양이, 옆모습 고양이, 자는 고양이, 점프하는 고양이…

로봇 행동도 똑같습니다. ``펜 뚜껑을 여는 방법''은 한 가지가 아니에요. 어느 손으로? 어느 각도로? 어느 속도로? 무수히 많은 valid trajectories가 있습니다.

[말줄임] 만약 단순한 regression으로 학습하면, 모델은 모든 valid trajectories의 \textbf{평균}을 출력하는 경향이 있어요. 그게 어떻게 보일까요? \textit{어색하고 비현실적인} 동작이 됩니다. 두 valid 동작의 평균이 valid한 동작이 아닌 거죠. 이걸 ``mode collapse''라고 합니다.

Flow matching은 이 문제를 \textbf{우회}합니다. 평균을 학습하지 않고, 분포 자체를 학습하니까요. 매번 \textit{여러 valid mode 중 하나}를 자연스럽게 샘플링합니다.

$\pi_0$의 절차를 정리하면---무작위 14×50 노이즈에서 출발해서, Action Expert가 ``이걸 어떻게 업데이트해야 더 그럴듯한 궤적이 되는지''를 예측하고, 더해주고, 다시 모델에 넣고, 또 업데이트하고. 이걸 \alert{10번 반복}하면 매끄러운 궤적이 나옵니다.
\end{script}
\end{frame}

% ============================================================
\begin{frame}{[19] 이미지 생성과 로봇 제어의 통합}
\begin{script}
잠시 한 발짝 떨어져서 봅시다. 우리가 방금 본 사실이 사실 굉장히 깊은 통찰을 담고 있어요.

\textbf{이미지 생성}에서: 텍스트와 노이즈 이미지를 입력으로 받아서, 자연 이미지의 manifold 위의 한 점을 출력합니다. ``cat''이라는 텍스트가 cat 이미지를 가리키도록요.

\textbf{로봇 제어 ($\pi_0$)}에서: 프롬프트, 이미지, 노이즈 궤적을 입력으로 받아서, 가능한 동작 manifold 위의 한 점을 출력합니다. ``uncap the pen''이라는 텍스트가 펜 따는 궤적을 가리키도록요.

\textbf{같은 수학, 다른 응용}이에요. 같은 flow matching 손실 함수, 같은 디노이징 절차, 같은 transformer 백본. 다른 건 데이터의 종류뿐입니다.

[말줄임] 이게 우연일까요? 아니라고 봅니다. \alert{매우 다른 응용 영역에서 같은 수학적 프레임워크가 작동한다}는 건, 그 프레임워크가 어떤 \textbf{심오한 통일성}을 포착하고 있다는 신호예요.

실제로 이 통찰이 후속 연구의 토대가 되었습니다. Diffusion Policy, RDT, Octo 같은 모델들이 모두 이 ``생성 모델 = 정책 학습'' 패러다임을 채택하고 있어요. 우리 lab의 ScaleMaster나 LT-Mem 같은 작업에서도 이 흐름을 어떻게 활용할지 고민해 볼 수 있습니다.
\end{script}
\end{frame}

% ============================================================
\section{Part 6. 두 모듈의 결합 --- KV cache 트릭}

\begin{frame}{[20] 문제 설정}
\begin{script}
이제 가장 영리한 부분으로 갑시다. 두 모듈을 어떻게 연결하느냐.

상황을 정리해 봅시다. Action Expert도 Gemma와 같은 구조예요. 18 layer × 8 head를 가집니다.

Action Expert 입장에서, 자기 query는 \textbf{무엇을 ``보고'' 싶을까요?} 네 가지가 있어요.

\textbf{첫째}, 다른 시간 단계의 행동들. 부드러운 궤적을 만들려면 ``이전 시간 단계에서 손이 어디 있었나''를 봐야 하니까요.

\textbf{둘째}, 프롬프트. ``무엇을 해야 하는가''를 알아야 합니다.

\textbf{셋째}, 이미지. ``펜이 어디 있는가''를 봐야죠.

\textbf{넷째}, 로봇 상태. ``지금 내 손이 어디 있는가''도 알아야 합니다.

[말줄임] 그런데 첫째와 넷째는 Action Expert 자기 안에 있어요. 그건 쉽습니다. 문제는 둘째와 셋째예요. 프롬프트와 이미지 정보는 \textbf{PaliGemma의 attention 안에 들어 있습니다}. 어떻게 거기에 접근하게 할까요?

이게 다음 슬라이드의 핵심 트릭, ``Key/Value 이어 붙이기''입니다.
\end{script}
\end{frame}

% ============================================================
\begin{frame}{[21] 해법: Key/Value를 이어 붙이기}
\begin{script}
해법은 단순하지만 영리합니다. 한 줄로 말하면, \alert{같은 layer index에서 두 모델의 K, V를 그냥 concatenate한다}는 거예요.

수식으로 봅시다. Action Expert의 layer $\ell$에서 attention을 계산할 때, key 행렬을 이렇게 만듭니다. Action Expert 자기 자신의 $K^\text{AE}_\ell$ 위에 PaliGemma의 $K^\text{Gemma}_\ell$을 쌓아서 \alert{$K_\text{total}$}을 구성하는 거예요. Value 행렬도 마찬가지로요.

크기를 보면 51 + 772 = \alert{823 keys}입니다. Action Expert의 query 51개가, 자기 51 토큰과 Gemma의 772 토큰을 \textbf{한 번의 attention 연산으로 동시에} 참조할 수 있게 되는 거죠.

\textit{왜 이게 가능할까요?} 두 모델이 \textbf{같은 width, 같은 layer 수, 같은 head 차원}을 갖기 때문입니다. Key/value 공간의 차원이 호환되니까 그냥 이어 붙일 수 있는 거예요.

[말줄임] 이게 바로 ``같은 아키텍처를 쓴다''는 설계 결정의 진짜 이유예요. 단순히 코드 재사용이 아니라, \textbf{이 KV concatenation을 가능하게 하기 위한 의도적 선택}이었던 겁니다.

다른 패러다임과 비교해 보세요. 만약 두 모델이 다른 아키텍처라면? Cross-attention layer를 추가하거나, projection layer를 둬서 차원을 맞추거나, 별도의 연결 메커니즘을 설계해야 합니다. 여기서는 그게 다 \emph{공짜}예요.
\end{script}
\end{frame}

% ============================================================
\begin{frame}{[22] 효율성: KV cache 재사용}
\begin{script}
KV concatenation의 또 다른 보너스가 있어요. 효율성입니다.

표준 LLM 추론에서 KV cache는 익숙한 개념이죠. 새 토큰이 들어올 때마다 이전 K, V를 재계산하지 않고 캐시에서 꺼내 쓰는 최적화예요.

$\pi_0$가 이걸 영리하게 활용합니다. 절차를 보면---

\textbf{1단계}: PaliGemma forward pass를 한 번 돌립니다. 이 과정에서 모든 layer의 $K^\text{Gemma}_\ell$과 $V^\text{Gemma}_\ell$을 캐시해 둡니다.

\textbf{2단계}: Flow matching은 \textbf{10번 반복}돼야 한다고 했죠? 매 반복마다 Action Expert가 새로 forward pass를 돕니다. 그런데---

\textbf{3단계}: 매 반복마다 \alert{같은 캐시를 재사용}합니다. 왜? \textbf{이미지와 프롬프트는 그동안 안 바뀌니까요}. PaliGemma forward를 다시 돌릴 필요가 없는 거예요.

[말줄임] 이 캐싱 덕분에 \textbf{10회 디노이징을 73 ms에 끝낼 수 있어요}. 이게 RTX 4090 한 장에서 실시간 제어가 가능해진 결정적인 이유입니다.

수치로 잠깐 감 잡아봅시다. 73 ms당 한 번 디노이징 사이클이 돌아간다는 건 약 \textbf{13 Hz의 제어 주기}예요. 인간의 반응 속도보다 빠릅니다. 이게 실제 로봇이 자연스럽게 움직이는 데 충분한 속도입니다.
\end{script}
\end{frame}

% ============================================================
\begin{frame}{[23] $\pi_0$의 한 step 전체 흐름}
\begin{script}
$\pi_0$ 한 사이클 전체를 한번에 정리해 봅시다.

\textbf{1단계}: 카메라 세 장 + 텍스트 프롬프트 + 현재 robot state를 수집합니다.

\textbf{2단계}: PaliGemma forward pass를 한 번 돌려서, 모든 layer의 K, V를 캐시합니다.

\textbf{3단계}: Action Expert를 초기화합니다. 1개의 state token + 50개의 noise token이 input이에요.

\textbf{4단계}: Flow matching 루프를 \textbf{10회} 돕니다. 각 회에서 self-attention과 (concatenated K/V를 통한) cross-attention을 동시에 수행하고, $\Delta a$를 계산해서 궤적을 업데이트합니다.

\textbf{5단계}: 최종 14×50 궤적을 출력합니다.

\textbf{6단계}: 로봇이 이 궤적의 \textbf{처음 몇 step만 실행}합니다. 끝까지 가는 게 아니에요. 왜냐하면 환경이 변할 수 있으니까요.

\textbf{7단계}: 새 카메라 이미지를 받아서 \textbf{1단계로 돌아갑니다}. 이걸 약 13 Hz로 반복하는 거죠.

[말줄임] 이게 ``model predictive control''의 일종으로 볼 수 있어요. \textbf{미래를 멀리 내다보고 계획하되, 일부만 실행하고 다시 계획하는} 패턴이죠. 고전 제어와 현대 학습 기반 정책의 흥미로운 결합입니다.
\end{script}
\end{frame}

% ============================================================
\section{Part 7. 비판적 시각과 미래}

\begin{frame}{[24] 데모를 신중히 읽기: 1995년 Ralph}
\begin{script}
지금까지 $\pi_0$의 놀라운 능력을 봤어요. 그런데 우리가 연구자로서 한 발짝 떨어져서 봐야 할 게 있습니다.

1995년 이야기를 잠깐 할게요. Carnegie Mellon에서 Ralph라는 자율주행 시스템을 만들었어요. 미국을 \textbf{98.2\% 자율}로 횡단했습니다. 1995년에요.

자, 이거 들으면 ``와, 그럼 자율주행차가 1995년에 다 됐나?'' 싶죠? \textit{전혀 아닙니다.} 30년이 지난 지금, 우리가 가진 자율주행차는 Ralph의 직계 후손이 아니에요. 센서 융합, HD map, learned planner, end-to-end transformer 같은 \textbf{전혀 다른 패러다임}의 조합이에요. 그리고 30년이 지나도 ``완전히 풀린'' 문제가 아닙니다.

[말줄임] $\pi_0$도 마찬가지예요. \textbf{인상적인 데모와 ``실제로 작동하는 시스템''은 다른 문제}입니다. 빨래 개기 데모가 인상적이라고 ``집안일 다 해주는 로봇''이 코앞에 있다고 결론 내리면 안 돼요.

대학원생으로서, 그리고 미래의 연구자로서, 데모를 \textbf{비판적으로 읽는 능력}이 정말 중요합니다. ``이 데모는 무엇을 보여주고, 무엇을 보여주지 \emph{않는가}?'' 이 질문을 항상 던져야 해요.
\end{script}
\end{frame}

% ============================================================
\begin{frame}{[25] 대안 패러다임: World Models}
\begin{script}
실제로 ``VLA가 정말 올바른 길인가?''에 대한 의문을 강하게 제기하는 학자들이 있어요. 가장 대표적인 사람이 \textbf{Yann LeCun}이에요.

LeCun은 AI 분야의 대선구자고, Meta의 수석 AI 과학자였어요. 그런데 최근에 Meta를 떠나서 \textbf{world model 중심의 새 벤처}를 시작했습니다. 그의 방향은 \alert{LLM을 백본으로 쓰지 않는 쪽}이에요. JEPA(Joint-Embedding Predictive Architecture) 계열이라고 합니다.

LeCun의 입장은 거침없어요. 인터뷰에서 ``JEPA가 결국 VLA를 추월할 거라 보시나요?'' 묻자, \textit{``물론이죠. VLA는 망했어요. 정말 잘 작동하지 않거든요''}라고 답했습니다.

[말줄임] LeCun이 무슨 말을 하고 있는지 핵심을 정리해 봅시다.

\textbf{VLA의 패러다임}: 인터넷 텍스트와 이미지로 사전학습 → 행동 데이터로 fine-tune. ``LLM이 곧 두뇌''라는 가정.

\textbf{World model의 패러다임}: 환경의 \textbf{예측 모델}을 latent space에서 학습. ``다음 순간이 어떻게 될지'' 시뮬레이션할 수 있는 능력 자체가 지능의 핵심이라는 가정. Planning은 학습된 world model 위에서 search나 optimization으로 합니다.

\textit{어느 쪽이 옳을까요?} 솔직히 아무도 모릅니다. 다만 두 패러다임이 강조하는 \textbf{근본적으로 다른 가정}을 이해하는 게 중요해요. 다음 강의에서 world model을 자세히 다루겠습니다.
\end{script}
\end{frame}

% ============================================================
\begin{frame}{[26] 토론거리}
\begin{script}
오늘 마무리는 토론으로 합시다. 다섯 가지 질문 던져드릴게요. 정답이 있는 질문이 아니에요. 다음 주까지 자기 생각 정리해 와 주세요.

\textbf{1.} RT2의 일반화(Taylor Swift)는 진짜 ``이해''인가, 아니면 통계적 보간일 뿐인가? 이걸 검증할 실험을 설계할 수 있을까?

\textbf{2.} Action Expert를 따로 두는 건 일종의 \textbf{귀납적 편향(inductive bias)}을 모델에 주입한 거예요. 이게 ``순수 end-to-end'' 철학과 충돌할까요, 아니면 보완할까요? 우리는 어디까지 architecture를 손봐야 할까요?

\textbf{3.} Flow matching이 mode collapse를 피한다고 했습니다. 그런데 \textbf{원치 않는 mode}---가령 안전하지 않은 동작---도 학습될 위험은 없을까요? 어떻게 막을 수 있을까요?

\textbf{4.} KV cache 공유는 ``같은 아키텍처''에 의존합니다. 그러나 미래의 로봇 시스템이 \textbf{이종 모듈}을 결합해야 한다면(예를 들어 SLAM 모듈 + LLM + 제어 모듈), 더 일반적인 방법은 무엇이 있을까요?

\textbf{5.} \textit{(가장 우리 lab에 해당하는 질문)} LeCun의 주장이 옳다면, classical SLAM과 state estimation은 미래 로봇 두뇌에서 어떤 자리를 차지할까요? 우리 APRL의 ``Classical-to-Modern Bridge'' 정체성을 어떻게 발전시켜야 할까요?

[말줄임] 다섯 번째 질문은 사실 우리 lab 정체성과 직결되는 부분이에요. ScaleMaster, LT-Mem, Mag4D-SLAM 모두 classical 기하학적 기초 위에서 modern learning을 결합하는 작업입니다. VLA 시대에서 이 결합의 의미를 어떻게 재정의할지가, 다음 5년의 우리 lab의 방향이 될 거라고 봐요.
\end{script}
\end{frame}

% ============================================================
\begin{frame}{[27] 핵심 정리}
\begin{script}
오늘 다룬 내용을 한 번에 정리합시다.

\textbf{1. VLA의 정의}: Vision + Language + Action을 단일 모델에서 통합하는 패러다임. 인터넷 사전학습의 지식을 행동과 연결하는 게 핵심 가치.

\textbf{2. $\pi_0$의 세 가지 설계 원칙}: 첫째, 백본 LLM(PaliGemma) 위에 별도 Action Expert를 둔다. 둘째, Action Expert는 flow matching으로 궤적을 생성한다. 셋째, KV cache 공유로 두 모듈을 매끄럽게 결합한다.

\textbf{3. 두 가지 깊은 통합}: 인터넷 지식 ↔ 물리적 행동의 통합과, 이미지 생성 수학 ↔ 로봇 제어 수학의 통합. 둘 다 ``서로 다른 영역이 같은 프레임워크로 연결된다''는 통일성을 보여줍니다.

\textbf{4. 한계와 대안}: World model 패러다임(LeCun, JEPA)은 ``LLM이 곧 두뇌''라는 VLA의 가정을 거부합니다. 어느 쪽이 옳은지는 아직 열려 있습니다.

\textbf{5. 연구자에게 주는 교훈}: 인상적인 데모와 실제로 작동하는 시스템 사이의 간극을 잊지 말 것. 1995년 Ralph가 우리에게 가르쳐 준 교훈입니다.

[말줄임] 마지막으로 한마디. 우리 lab은 \textbf{classical과 modern을 잇는 다리}라는 정체성을 갖고 있어요. VLA는 modern의 최전선이지만, 이게 잘 작동하려면 결국 \textbf{state estimation, geometric reasoning} 같은 classical 기초가 필요합니다. 정밀한 manipulation에는 결국 정확한 6DoF state가 필요한 것처럼요.
\end{script}
\end{frame}

% ============================================================
\begin{frame}{[28] 참고 문헌}
\begin{script}
오늘 강의 자료의 원천을 다시 명시합니다.

\textbf{원본 영상}: ``Inside the World's Smartest Robot Brain [VLA]'', \textbf{Welch Labs}.\\
URL: \href{https://www.youtube.com/watch?v=2mrGMMmrVNE}{youtube.com/watch?v=2mrGMMmrVNE}

영상이 매우 잘 만들어져 있으니 직접 보시는 걸 강력 추천합니다. 시각적 설명이 뛰어나요.

\textbf{핵심 논문 시작점}:
\begin{itemize}
  \footnotesize
  \item Ahn et al. (2022), SayCan
  \item Brohan et al. (2022), RT-1
  \item Driess et al. (2023), PaLM-E
  \item Brohan et al. (2023), RT-2
  \item Black et al. (2024), $\pi_0$
  \item Lipman et al. (2023), Flow Matching
  \item LeCun (2022), JEPA position paper
\end{itemize}
\end{script}
\end{frame}

% ============================================================
\begin{frame}{[29] 다음 주 예고와 마무리}
\begin{script}
다음 주에는 \textbf{World Model과 JEPA 계열}을 깊이 다뤄볼 거예요.

오늘이 ``LLM 기반 로봇 두뇌'' 편이었다면, 다음 주는 ``LLM에 의존하지 않는 로봇 두뇌'' 편입니다. 두 패러다임을 직접 대비시켜 볼 거예요.

\textbf{과제}: 오늘 토론거리 다섯 개에 대한 자기 생각을 짧게라도 정리해서 다음 주에 가져와 주세요. 특히 다섯 번째---``classical SLAM의 미래 자리''---는 우리 lab의 방향과 직결되니까 신중히 생각해 보시기 바랍니다.

\textbf{권장 예습}: LeCun의 2022년 JEPA position paper(``A Path Towards Autonomous Machine Intelligence''). 길지만 도입부와 결론만 봐도 다음 주 흐름을 잡을 수 있어요.

[말줄임] 오늘 수고하셨습니다. 질문 있으면 받겠습니다.
\end{script}
\end{frame}

\end{document}
